\documentclass[../../main.tex]{subfiles}
\graphicspath{{../../image/}{../../figure/}}
\begin{document}

\section{材料标架无差异原理}\label{section:ch02-材料标架无差异}

体力和热量供给通常由外部因素诱导并预先给定，而内能、应力、熵和热通量这些场则由热力学过程确定。运动可在其使物体变形的程度上影响这些场：刚体运动不改变质点之间的距离，因而不应对内能、温度或参考热通量产生影响；但它应以如下方式影响应力张量，使得从附着于运动物体的标架观察时，所得应力向量看起来是固定的。这一要求由\emph{材料标架无差异基本原理}所公设，现在我们精确陈述之。

考察物体的任意两个热力学过程 $(\chi,s)$ 和 $(\chi^{\#},s^{\#})$，使得熵场重合 $s^{\#}=s$，而运动相差一个刚性的（依赖于时间的）旋转\footnote{这一设定的另一种等价实现是：将一个单一的热力学过程设想为由附着于各自坐标系标架的两个观察者监测，而这两个标架彼此相对旋转。采用该途径时，某些作者允许除正常旋转外还包含反射。}：
\begin{align}
\chi^{\#}(x,t)=O(t)\chi(x,t),\qquad x\in\mathscr{B},\ t\in(t_{1},t_{2}),
\label{align:ch02-2-4-1}
\end{align}
\begin{align}
O^{\top}(t)O(t)=O(t)O^{\top}(t)=I,\qquad\det O(t)=1,\quad t\in(t_{1},t_{2}).
\label{align:ch02-2-4-2}
\end{align}
注意两个过程 $(\chi,s)$、$(\chi^{\#},s^{\#})$ 的变形梯度 $F,F^{\#}$、空间速度梯度 $L,L^{\#}$ 和拉伸张量 $D,D^{\#}$（参见\eqref{align:ch02-2-1-8}、\eqref{align:ch02-2-1-9}）由下式联系：
\begin{align}
F^{\#}=OF,\qquad L^{\#}=OLO^{\top}+\dot{O}O^{\top},\qquad D^{\#}=ODO^{\top}.
\label{align:ch02-2-4-3}
\end{align}

设 $(\varepsilon,S,\theta,Q)$ 和 $(\varepsilon^{\#},S^{\#},\theta^{\#},Q^{\#})$ 分别为与过程 $(\chi,s)$ 和 $(\chi^{\#},s^{\#})$ 相伴随的内能、Piola--Kirchhoff 应力、温度和参考热通量场。材料标架无差异原理公设：
\begin{align}
\varepsilon^{\#}=\varepsilon,\qquad S^{\#}=OS,\qquad\theta^{\#}=\theta,\qquad Q^{\#}=Q.
\label{align:ch02-2-4-4}
\end{align}
由\eqref{align:ch02-2-4-4}、\eqref{align:ch02-2-3-17} 和 \eqref{align:ch02-2-4-3} 推得，两个过程的参考与空间温度梯度 $G,G^{\#}$ 和 $g,g^{\#}$ 由下式联系：
\begin{align}
G^{\#}=G,\qquad g^{\#}=Og.
\label{align:ch02-2-4-5}
\end{align}
此外，由\eqref{align:ch02-2-3-6}、\eqref{align:ch02-2-3-12} 和 \eqref{align:ch02-2-4-3}，我们推出两个过程的 Cauchy 应力张量 $T,T^{\#}$ 与空间热通量向量 $q,q^{\#}$ 之间的如下关系：
\begin{align}
T^{\#}=OTO^{\top},\qquad q^{\#}=Oq.
\label{align:ch02-2-4-6}
\end{align}

材料标架无差异原理应反映在连续介质的本构关系中，无论材料响应的性质如何。说明性例子将在随后两节中考察。

\end{document}
